% Technical note --- rationale for the OCR engine selection
% Build:  pdflatex technical_note_ocr_en.tex  (twice, for the table of contents)
\documentclass[11pt,a4paper]{article}

\usepackage[utf8]{inputenc}
\usepackage[T1]{fontenc}
\usepackage{lmodern}   % vector T1 fonts (avoids bitmap ec fonts in the output)
\usepackage[english]{babel}
\usepackage[a4paper,margin=2.5cm]{geometry}
\usepackage{graphicx}
\usepackage{tabularx}
\usepackage{enumitem}
\usepackage{array}
\usepackage{float}     % [H] keeps every table where it is written

\newcolumntype{L}[1]{>{\raggedright\arraybackslash}p{#1}}

% booktabs-style rules, defined locally so the note builds on a bare LaTeX
% installation (booktabs is not part of every MiKTeX setup).
\newcommand{\toprule}{\noalign{\hrule height 0.9pt}\noalign{\vskip 3pt}}
\newcommand{\midrule}{\noalign{\vskip 3pt}\noalign{\hrule height 0.4pt}\noalign{\vskip 3pt}}
\newcommand{\bottomrule}{\noalign{\vskip 3pt}\noalign{\hrule height 0.9pt}}

\setlist[itemize]{itemsep=2pt,topsep=4pt}
\setlist[enumerate]{itemsep=2pt,topsep=4pt}

\begin{document}

% ---------------------------------------------------------------- title page
\begin{titlepage}
\centering
\vspace*{5cm}
\includegraphics[width=4cm]{../ey.png}

\vspace{2.5cm}
{\LARGE\bfseries Intelligent OCR for Scanned Documents\par}

\vspace{0.8cm}
{\Large French \quad --- \quad English \quad --- \quad Arabic\par}

\vspace{0.8cm}
\rule{0.45\textwidth}{0.4pt}

\vspace{0.8cm}
{\large Technical note: rationale for the OCR engine selection\par}

\vfill
{\large July 2026\par}
\vspace{1cm}
\end{titlepage}

\tableofcontents
\newpage

% ---------------------------------------------------------------- 1
\section{Context and Objectives}

The project addresses the need to digitise and automatically process
\textbf{scanned} documents: administrative forms, reports and handwritten
material, written in \textbf{French}, \textbf{English} or \textbf{Arabic}.

Classical OCR engines (Tesseract, for example) detect \emph{character shapes}
with no understanding of context, layout or language. They degrade sharply on
noisy scans, mixed scripts, tables and handwriting --- precisely the cases
encountered here.

By contrast, a \textbf{Vision-Language Model} (VLM) reads an image and produces
text within a \emph{single} neural network. It is this joint understanding of
image and language that makes difficult documents tractable.

\paragraph{Objectives.}
\begin{itemize}
  \item Faithfully transcribe multilingual scanned documents (FR / EN / AR,
        including right-to-left script).
  \item Preserve the document \textbf{structure} (headings, tables, checkboxes),
        not merely the raw text.
  \item Guarantee \textbf{confidentiality}: the reference document used for
        testing is a strictly confidential safeguarding form containing names
        and medical details. No document may leave the machine.
  \item Run on the hardware actually available: a laptop fitted with an
        \textbf{NVIDIA RTX 4050 (6\,GB VRAM)}.
  \item Incur no recurring cost: no per-page billing, no quota.
  \item Deliver the capability as a usable product, not only as a script: a
        multi-user web platform with authentication and a persistent
        per-user document history.
\end{itemize}

% ---------------------------------------------------------------- 2
\section{Functional Scope}

\paragraph{In scope.}
\begin{itemize}
  \item \textbf{Inputs}: images (PNG / JPG / WEBP) and scanned PDFs
        (multi-page, rasterised locally).
  \item \textbf{Languages}: French, English, Arabic --- including mixed FR/AR
        documents and non-Latin right-to-left scripts.
  \item \textbf{Content handled}: dense text, tables, structured forms,
        checkboxes and tick-boxes, handwriting.
  \item \textbf{Output}: structured text in \textbf{Markdown}, preserving the
        hierarchy and tables of the source document, with checkbox states
        reported explicitly.
  \item \textbf{Execution modes}: local engine by default, second opinion
        (quality assurance) on a larger model, optional cloud engine.
  \item \textbf{Quality control without ground truth}: recall verification,
        two-model completeness comparison, and format validation of
        constrained fields.
  \item \textbf{Interfaces}: a full-stack web platform (primary deliverable),
        plus a Streamlit prototype and a command-line interface over the same
        OCR core.
  \item Full \textbf{offline operation}.
\end{itemize}

\paragraph{Out of scope.}
\begin{itemize}
  \item Cloud processing as the default: ruled out for data-governance reasons.
  \item Business validation of the extracted content: the system transcribes;
        it neither interprets nor validates medical information.
  \item High-throughput batch processing: the target is the individual
        document, at roughly 30 seconds per page.
  \item Public deployment: there is no user signup flow by design; accounts are
        seeded or created directly in the database.
\end{itemize}

% ---------------------------------------------------------------- 3
\section{Overall Architecture}

The architecture rests on one guiding principle: \textbf{local-first}. The
nominal engine runs entirely on the user's machine; the cloud is only an
optional alternative.

\paragraph{Processing chain.}
\begin{enumerate}
  \item \textbf{Ingestion}: loading the image or scan. PDFs are rasterised
        locally, one image per page.
  \item \textbf{Preparation}: image normalisation (denoising, deskew, contrast)
        and construction of the transcription prompt.
  \item \textbf{VLM inference}: reading the image and generating the text with
        the vision-language model, served locally.
  \item \textbf{Structuring}: Markdown formatting of the output (headings,
        tables, lists), language detection and right-to-left handling.
  \item \textbf{Quality control (optional)}: recall verification pass, second
        opinion on a larger model, format validation of constrained fields.
  \item \textbf{Delivery}: hand-off of the transcribed document --- displayed,
        downloadable, and stored in the user's history.
\end{enumerate}

\paragraph{Engines involved.}
Four engines sit behind a single abstraction, so any of them can be swapped
without touching the processing chain.

\begin{table}[H]
\centering
\begin{tabular}{@{}L{3.2cm}L{5.6cm}L{5.2cm}@{}}
\toprule
\textbf{Role} & \textbf{Model} & \textbf{Where it runs} \\
\midrule
Default engine & Qwen3-VL 4B (instruct) --- \texttt{qwen3-vl:4b-instruct} &
  Local, via Ollama, on the GPU \\
Second opinion (QA) & Qwen3-VL 8B (instruct) --- \texttt{qwen3-vl:8b-instruct} &
  Local, via Ollama \\
Classical baseline & Tesseract (\texttt{fra+eng+ara}) &
  Local, optional --- benchmarking and hybrid mode \\
Optional cloud & Llama 4 Scout (vision) & Groq API \\
\bottomrule
\end{tabular}
\caption{Division of roles between the engines}
\end{table}

% ---------------------------------------------------------------- 4
\section{Detailed Technical Architecture}

\subsection{Image preprocessing}

Preprocessing runs on \textbf{OpenCV}, with a Pillow-only fallback when OpenCV
is unavailable: grayscale conversion, non-local-means denoising, deskew
(estimation of the dominant text angle, applied only between $0.3^\circ$ and
$15^\circ$ to avoid acting on a misdetection), and CLAHE contrast enhancement.
Adaptive binarisation is applied \emph{only} on the Tesseract path: the vision
model performs measurably better on contrast-enhanced grayscale than on a
hard-thresholded image.

Images are normalised to the model's useful resolution --- long edge capped at
2576\,px, small scans upscaled --- and encoded as JPEG to keep the payload
within the limits of the vision APIs.

\subsection{Local inference service}

The models are served by \textbf{Ollama}, which exposes a local API and manages
loading the weights onto the GPU. No outbound network call is required for
nominal operation.

\subsection{Memory constraint and GPU residency}

The RTX 4050 provides \textbf{6\,GB of VRAM}. The 4B model occupies about
\textbf{3.3\,GB} and therefore stays \emph{fully resident} on the GPU. Larger
models spill over to the CPU, which severely degrades response times.

\subsection{The \texttt{-instruct} suffix is mandatory}

A decisive technical point: the base \texttt{qwen3-vl} tags run in
\emph{thinking} mode. They spend the entire token budget on internal reasoning
and return an \textbf{empty transcription}. Only the \texttt{-instruct} variant
transcribes the image content directly. The engine detects this failure mode
explicitly and reports it rather than returning an empty result.

\subsection{Context-window tuning}

Setting the \texttt{num\_ctx} parameter to \textbf{8192 tripled the speed}
compared with the naive initial setting, bringing processing down to roughly
\textbf{30 seconds per page} on the target hardware. The value is a genuine
optimum rather than a maximum: each page is a separate call, so the window only
needs to hold one page (roughly 4300 image tokens plus the transcription).
Raising it further wastes VRAM on the KV cache and forces model layers onto the
CPU --- measured about three times slower per page, for no accuracy gain.

\subsection{Multi-page PDF handling}

Neither engine accepts PDF input natively. Pages are rasterised locally with
\textbf{PyMuPDF} at 300\,DPI (no external Poppler dependency) and transcribed
one at a time, the results being joined with a page separator. The resolution
matters: fine print --- form placeholders, footnotes --- sits right at the
model's legibility threshold, and 200--250\,DPI drops it where 300\,DPI reads
it reliably. Pages are processed sequentially: on a single GPU that is the fast
path, since concurrent calls merely contend for the same VRAM.

\subsection{Prompt design}

Three system prompts drive the model, all sharing a strict
``transcribe, never interpret'' contract --- no translation, no summarising, no
grammar correction:
\begin{itemize}
  \item \textbf{Extraction}: reproduce every character faithfully, including
        accents and Arabic diacritics; transcribe small print, footnotes,
        headers, stamps and form helper text; report the state of every
        checkbox next to its label; preserve reading order, right-to-left for
        Arabic passages; render structure as Markdown, tables as Markdown
        tables; mark genuinely illegible regions rather than guessing.
  \item \textbf{Proofreading}: used by the verification pass (Section 6.3).
  \item \textbf{Post-correction}: used by the hybrid mode to repair classical
        OCR output --- character confusions, broken words, missing diacritics
        --- without paraphrasing.
\end{itemize}

\subsection{Cloud path (optional)}

A much larger remote engine --- \textbf{Llama 4 Scout} (17B, vision) via the
\textbf{Groq} API --- is integrated as a benchmark and an alternative for cases
willing to trade confidentiality for maximum accuracy. It shares the same
prompts and the same processing chain as the local engine, so local and cloud
results are directly comparable. It is selected per request and requires an API
key; the local engine remains the default and needs none.

% ---------------------------------------------------------------- 5
\section{Recognition Modes}

Three modes run over the same chain, selectable per request in every interface:

\begin{enumerate}
  \item \textbf{VLM vision} (recommended, default) --- the preprocessed page is
        read directly by the vision-language model. Best on degraded scans,
        complex layouts, tables and Arabic script.
  \item \textbf{Hybrid} --- Tesseract produces a fast raw transcription, then
        the VLM repairs character confusions, broken words and accent errors
        without paraphrasing.
  \item \textbf{Tesseract only} --- the classical baseline, no VLM involved.
        Useful for benchmarking the VLM approach against a traditional one.
\end{enumerate}

Modes 2 and 3 require the Tesseract binary with the \texttt{fra}, \texttt{eng}
and \texttt{ara} language packs. Tesseract is auto-detected: when it is absent
--- the normal configuration --- both modes are simply hidden from the
interfaces, and only VLM vision is offered.

% ---------------------------------------------------------------- 6
\section{Quality Control and Validation}

\subsection{Measurement protocol}

The comparison was run on the most demanding case in the corpus: a
\textbf{dense handwritten form} containing tables and checkboxes. The chosen
criterion is transcription \textbf{completeness}, supplemented by processing
time. Where a ground truth is available, \textbf{CER} and \textbf{WER}
(character and word error rate, Levenshtein-based, with Unicode NFC
normalisation --- required for consistent Arabic comparison) are computed
instead.

\begin{table}[H]
\centering
\begin{tabular}{@{}L{4.2cm}ccL{5.4cm}@{}}
\toprule
\textbf{Model} & \textbf{Completeness} & \textbf{Duration} & \textbf{Observation} \\
\midrule
Qwen3-VL 4B instruct & 10/10 & 59\,s & Fully resident on the GPU \\
Qwen3-VL 8B instruct & 9/10 & 135\,s & Spilled to CPU (6\,GB VRAM insufficient) \\
\bottomrule
\end{tabular}
\caption{Results measured on the reference handwritten form}
\end{table}

The result is counter-intuitive but unambiguous: on this machine the 4B model
is both \textbf{faster and more complete} than the 8B. The choice is therefore
not assumed, it is \textbf{measured}.

\subsection{The dominant failure mode is omission}

On a dense form, a vision model fails far more often by \emph{skipping} text
than by inventing it: faint print, marginal notes, table cells and checkboxes
get silently dropped. Simply reading the page twice does not help --- the model
tends to skip the same things twice. Three independent tools address this from
different angles.

\subsection{Recall verification (two-pass)}

The verification option runs a two-pass loop. The first pass is a normal
transcription. The second pass gives the model \emph{the same image plus its own
draft}, under a prompt whose only job is to find what the draft missed and fix
what it got wrong, while being forbidden to remove correct text, translate,
reorder or summarise.

Handing the model its own output as a checklist turns an open-ended
transcription task into a targeted search, which is what recovers the skipped
spans. The cost is exactly \textbf{twice the model calls} --- so about twice the
time and twice the tokens, per page on a PDF.

\subsection{Two-model completeness comparison}

No OCR engine --- local or cloud --- can \emph{guarantee} it never drops a word;
accuracy is always a measured rate, not an absolute. The honest way to push
towards completeness is to read each page with \textbf{two independent models}
and review wherever they disagree: agreement means high confidence,
disagreement marks a spot worth a human glance.

The comparison mode does exactly this, fully offline, using the 4B and 8B local
models. It reports a word-level agreement percentage and every disagreement
span, flagging text one model skipped and the other caught. In practice the two
models have \emph{different} blind spots --- the 4B is stronger on some
checkboxes and fine print, the 8B catches some headers the 4B drops --- so the
diff reliably surfaces real omissions that a single pass would hide.

\subsection{Format validation without ground truth}

Digits and proper nouns are where a vision model fails first: they carry no
linguistic context to fall back on, so a misread \texttt{B}~$\rightarrow$~\texttt{8}
produces something perfectly plausible. On a form, those are exactly the fields
that matter --- a wrong postcode or phone number is worse than a wrong
adjective.

Those fields are \emph{format-constrained}, so a share of that error class is
catchable \textbf{with no ground truth at all}. Three checks run over the
transcription:

\begin{table}[H]
\centering
\begin{tabular}{@{}L{2.8cm}L{5.4cm}L{5.6cm}@{}}
\toprule
\textbf{Check} & \textbf{What it catches} & \textbf{Example} \\
\midrule
Structural & The token violates the field's grammar &
  A UK postcode cannot hold a digit in its final two positions \\
Consistency & The same entity transcribed two ways &
  Two spellings of one postcode, both legal --- so one is misread \\
Plausibility & Wrong length &
  A 9- or 10-digit UK phone number means a dropped digit \\
\bottomrule
\end{tabular}
\caption{Format checks applied to constrained fields}
\end{table}

When exactly \emph{one} OCR confusion substitution
(\texttt{8}~$\rightarrow$~\texttt{B}, \texttt{5}~$\rightarrow$~\texttt{S},
\texttt{2}~$\rightarrow$~\texttt{Z}, \dots) repairs a malformed token, the fix
is unambiguous and is printed as a suggestion; ambiguous repairs are reported
without one. The rules are currently UK-specific, because that is what the test
corpus contains.

This validates \textbf{form, not truth}: a well-formed but wrong value --- a
valid postcode for the wrong town --- passes clean. A clean report means ``no
detectable format errors'', never ``no errors''.

\subsection{Known limitations}
\begin{itemize}
  \item \textbf{Hardware dependency.} Performance is tied to the presence of a
        GPU and its VRAM; on a machine without a GPU, processing times become
        prohibitive.
  \item \textbf{Scope of the measurement.} The figures above come from one
        reference document on one machine; they indicate a trend, not a
        universal guarantee.
  \item \textbf{Fidelity not guaranteed at 100\%.} Given that the data is
        personal and medical, \textbf{human proofreading} of the transcription
        remains necessary. None of the three quality tools replaces it.
  \item \textbf{Arabic handwriting.} It remains the hardest case in the chain
        and accounts for most of the residual errors.
  \item \textbf{Third-party dependency.} The cloud path depends on a provider's
        model catalogue, which changes over time; the model identifier is
        therefore configurable rather than hard-coded.
  \item \textbf{Validation rules are locale-specific.} The format checks encode
        UK postcode and phone grammar and would need extending for other
        locales.
\end{itemize}

% ---------------------------------------------------------------- 7
\section{Technology Choices}

\subsection{Why a VLM rather than classical OCR}

Classical OCR recognises glyphs; a VLM \emph{understands} the page. That
understanding of context, layout and language is what makes it possible to
handle degraded scans, mixed scripts, tables and handwriting. Tesseract is
nonetheless retained as an optional baseline, so the claim can be verified by
measurement rather than asserted.

\subsection{Why the Qwen3-VL family}
\begin{itemize}
  \item \textbf{Strong multilingual performance, Arabic included.} The project
        targets French, English and Arabic (right-to-left). Qwen models are
        trained heavily on multilingual and Arabic data, and rank among the best
        open VLMs on non-Latin scripts --- confirmed on the mixed FR/AR
        documents in the corpus.
  \item \textbf{Built for document tasks.} They handle dense text, tables and
        structured forms, and produce structured \textbf{Markdown} output:
        exactly what the reference form required, with its tables and
        checkboxes.
  \item \textbf{Open weights, local execution.} Usable through Ollama with no
        licence fee.
\end{itemize}

\subsection{Why the 4B size, and why the \emph{instruct} variant}
\begin{itemize}
  \item \textbf{4B fits the hardware.} About 3.3\,GB, meaning full residency
        within the 6\,GB of VRAM on the RTX 4050.
  \item \textbf{Measured, not assumed.} 10/10 in 59\,s for the 4B against 9/10
        in 135\,s for the 8B (see Section 6.1).
  \item \textbf{The \texttt{-instruct} tag is imperative.} Without it, the model
        reasons instead of transcribing and returns an empty output.
\end{itemize}

\subsection{Why local (Ollama) is the default engine}
\begin{itemize}
  \item \textbf{Confidentiality.} The test document is a strictly confidential
        safeguarding form with names and medical details. Local processing means
        \emph{no document ever leaves the machine}: the decisive factor for data
        governance.
  \item \textbf{No cost, no limits, offline.} No per-page API fees, no rate
        limits, no need for an internet connection.
  \item \textbf{Fast enough on real hardware.} Roughly 30 seconds per page after
        tuning the context window.
\end{itemize}

\subsection{Why Groq is kept as an option}

This remote engine serves as a \textbf{benchmark} and an \textbf{alternative}
towards a substantially larger model, for cases that would accept trading
confidentiality for maximum accuracy. Keeping both engines behind one
abstraction also means the local-first choice is a measured decision that can be
re-tested at any time, rather than a constraint of the design.

\subsection{Why this delivery stack}

The OCR core is deliberately interface-agnostic: the web backend, the prototype
and the command line all call the same pipeline object. The web platform is
built on \textbf{React} (Vite) for the interface, \textbf{FastAPI} for the API
--- native multipart upload handling and a natural fit for the Python OCR core
--- and \textbf{MongoDB} for storage, whose document model matches transcription
records with heterogeneous metadata better than a relational schema would.
Authentication uses JWT with PBKDF2-SHA256 password hashing.

% ---------------------------------------------------------------- 8
\section{Delivered Interfaces}

Three interfaces share the same OCR core:

\begin{itemize}
  \item \textbf{Web platform} (primary deliverable) --- an EY-branded
        single-page application: authentication, drag-and-drop upload of images
        and PDFs, engine and language selection, live progress indication,
        right-to-left rendering for Arabic-heavy results, per-run metrics
        (engine, languages, duration, token usage), copy and download, and a
        permanent searchable per-user history persisted in MongoDB.
  \item \textbf{Streamlit prototype} --- a single-file interface for quick
        one-off tests, showing the original and preprocessed images side by
        side and offering CER/WER evaluation against a pasted ground truth.
  \item \textbf{Command-line interface} --- for scripting and benchmarking; it
        additionally exposes the three quality-control tools of Section 6
        (verification pass, two-model comparison, format validation).
\end{itemize}

% ---------------------------------------------------------------- 9
\section{Project Structure}

The breakdown below reflects the separation between the OCR engines, the
processing chain, evaluation and delivery.

\begin{table}[H]
\centering
\begin{tabular}{@{}L{3.4cm}L{10.6cm}@{}}
\toprule
\textbf{Component} & \textbf{Responsibility} \\
\midrule
Engine layer & Common abstraction over the OCR engines: local implementation
  (Ollama), cloud implementation (Groq), and the optional classical baseline
  (Tesseract). \\
Preprocessing & Image normalisation: denoising, deskew, contrast, resolution
  capping, PDF rasterisation. \\
Processing chain & Ingestion, preparation, engine invocation, language
  detection, Markdown structuring. \\
Configuration & Model selection, inference parameters (including
  \texttt{num\_ctx}), choice of default engine and second-opinion engine. \\
Quality control & Verification pass, two-model comparison, format validation,
  CER/WER evaluation. \\
Web backend & REST API, authentication, per-user document persistence. \\
Web frontend & Single-page application: upload, engine selection, result
  viewing, history. \\
Outputs & Markdown transcriptions produced, kept locally. \\
\bottomrule
\end{tabular}
\caption{Functional breakdown of the project}
\end{table}

This breakdown makes it possible, in particular, to swap engines --- local or
cloud --- without touching the processing chain, and to add an interface
without touching either.

% ---------------------------------------------------------------- 10
\section{Performance}

Measured on the target laptop (RTX 4050, 6\,GB VRAM) with the default local
engine at 300\,DPI and \texttt{num\_ctx}~=~8192:

\begin{table}[H]
\centering
\begin{tabular}{@{}L{7.4cm}L{6.6cm}@{}}
\toprule
\textbf{Workload} & \textbf{Time} \\
\midrule
Printed page & 25--29\,s \\
Dense handwritten page & 32--47\,s \\
Three-page printed PDF & 86\,s (measured) \\
Ten-page form & about 5 minutes \\
First page after the service is idle & +10\,s, one-time (model loads into VRAM) \\
Verification pass & about 2$\times$ (two model calls per page) \\
Two-model comparison & 2.5--3$\times$ (runs both the 4B and the 8B) \\
Format validation & negligible (no model call) \\
\bottomrule
\end{tabular}
\caption{Measured processing times}
\end{table}

Rule of thumb: \textbf{time $\approx$ pages $\times$ 30\,s}, plus a one-time
10\,s on the first scan after a reboot. Two parameters trade speed against
completeness: the context window --- the dominant lever, see Section 4.5 ---
and the PDF rasterisation resolution, where lowering 300 to 250\,DPI saves about
20\,s per page but can drop the finest print. Both defaults deliberately favour
completeness over speed.

% ---------------------------------------------------------------- 11
\section{Development Planning}

\begin{table}[H]
\centering
\begin{tabular}{@{}cL{3.6cm}L{9.0cm}@{}}
\toprule
\textbf{Phase} & \textbf{Title} & \textbf{Content} \\
\midrule
1 & Study and framing & Analysis of the multilingual requirement, assessment of
  the limits of classical OCR, identification of the confidentiality and
  hardware constraints. \\
2 & Model selection & Review of open VLMs, selection of the Qwen3-VL family for
  its multilingual and document-handling performance. \\
3 & Local prototyping & Bringing Ollama into service, first transcription
  trials, identification of the \emph{thinking}-mode problem and switch to the
  \texttt{-instruct} variants. \\
4 & Optimisation & Context-window tuning (\texttt{num\_ctx}~=~8192), speed
  tripled, verification of GPU residency. \\
5 & Comparative evaluation & 4B versus 8B on the reference handwritten form;
  measurement of completeness and timings. \\
6 & Cloud integration & Addition of the Groq / Llama 4 Scout engine as a
  benchmark and alternative behind the shared engine abstraction. \\
7 & Multi-page and preprocessing & PDF rasterisation at 300\,DPI, page-by-page
  transcription, OpenCV preprocessing chain. \\
8 & Quality control & Verification pass, two-model completeness comparison and
  format validation of constrained fields. \\
9 & Web platform & REST API, authentication and per-user persistence; branded
  single-page interface with upload, results and history. \\
10 & Consolidation & Fixing the default engine, setting up the second opinion,
  documenting the known limitations. \\
\bottomrule
\end{tabular}
\caption{Development phases}
\end{table}

% ---------------------------------------------------------------- 12
\section{Conclusion}

The system uses \textbf{Qwen3-VL-4B-Instruct}, an open vision-language model
served locally via \textbf{Ollama}, as its default OCR engine.

The choice rests on three converging arguments: its multilingual performance,
notably in Arabic; its ability to handle documents, tables and handwriting; and
--- critically --- the fact that it runs \textbf{fully offline on a 6\,GB laptop
GPU}, keeping confidential documents on the device.

The 4B size was selected over the 8B after \textbf{direct measurement}, which
showed it to be both faster and more complete on the target hardware. A cloud
engine (Groq / Llama 4 Scout) is retained behind the same abstraction as a
higher-capacity alternative and as a benchmark against which the local-first
choice can be re-verified.

Around that engine, the project delivers a complete chain rather than a single
call: image preprocessing, page-by-page PDF handling, structured Markdown
output, three quality-control tools that address the dominant failure mode ---
omission --- from different angles, and a multi-user web platform through which
the whole thing is actually usable. What none of it removes is the need for
human proofreading of a critical document; the system is built to make that
review short, not to make it unnecessary.

\end{document}
